\documentclass[11pt,a4paper,oneside]{book}

% ==================================================================
%  Ultime Team Manager - Login paso a paso (estilo libro)
%  Compilar con:  tectonic login-paso-a-paso.tex
% ==================================================================

\usepackage[T1]{fontenc}
\usepackage[utf8]{inputenc}
\usepackage[spanish,es-noshorthands]{babel}
\usepackage[a4paper,margin=2.3cm]{geometry}
\usepackage{xcolor}
\usepackage{listings}
\usepackage{tcolorbox}
\usepackage{titlesec}
\usepackage{fancyhdr}
\usepackage{enumitem}
\usepackage{hyperref}

\tcbuselibrary{skins,breakable}

% ---- Paleta ------------------------------------------------------
\definecolor{brand}{HTML}{0B7A4B}
\definecolor{brandDark}{HTML}{064D2F}
\definecolor{accent}{HTML}{1565C0}
\definecolor{gold}{HTML}{C9A227}
\definecolor{codebg}{HTML}{F7F8FA}
\definecolor{codeframe}{HTML}{D7DCE1}
\definecolor{kw}{HTML}{9C27B0}
\definecolor{type}{HTML}{0B7A4B}
\definecolor{cmt}{HTML}{7A8B99}
\definecolor{str}{HTML}{C2410C}
\definecolor{numgray}{HTML}{9AA5B1}
\definecolor{midgray}{HTML}{5B6572}

\hypersetup{
  colorlinks=true, linkcolor=brandDark, urlcolor=accent,
  pdftitle={Ultime Team Manager - Login paso a paso},
  pdfauthor={Sebastian Ramirez, Joaquin Parraud}
}

% ---- Encabezado / pie -------------------------------------------
\pagestyle{fancy}
\fancyhf{}
\renewcommand{\headrulewidth}{0.4pt}
\lhead{\small\textcolor{midgray}{Ultime Team Manager}}
\rhead{\small\textcolor{midgray}{\leftmark}}
\cfoot{\small\textcolor{midgray}{\thepage}}

% ---- "Fase" en vez de "Capitulo" --------------------------------
\addto\captionsspanish{\renewcommand{\chaptername}{Fase}}
\titleformat{\chapter}[display]
  {\normalfont\bfseries\color{brandDark}}
  {\filright\Large\chaptertitlename\ \thechapter}{6pt}{\Huge\color{brand}}
\titlespacing*{\chapter}{0pt}{10pt}{24pt}

\titleformat{\section}
  {\normalfont\large\bfseries\color{accent}}{\thesection}{0.6em}{}

% ---- Estilo de codigo (Dart) ------------------------------------
\lstdefinelanguage{Dart}{
  morekeywords={abstract,as,assert,async,await,break,case,catch,class,const,
    continue,covariant,default,deferred,do,dynamic,else,enum,export,extends,
    extension,external,factory,false,final,finally,for,Function,get,hide,if,
    implements,import,in,interface,is,late,library,mixin,new,null,on,operator,
    part,required,rethrow,return,sealed,set,show,static,super,switch,sync,this,
    throw,true,try,typedef,var,void,while,with,yield},
  morekeywords=[2]{String,int,double,bool,Future,Widget,List,Map,BuildContext,
    Object,Exception,override,Duration,GlobalKey,FormState,TextEditingController,
    Notifier,Provider,NotifierProvider,ConsumerWidget,ConsumerStatefulWidget,
    ConsumerState,WidgetRef,Ref,GoRouter,GoRoute,ValueNotifier},
  sensitive=true,
  morecomment=[l]{//},
  morecomment=[s]{/*}{*/},
  morestring=[b]',
  morestring=[b]",
}

\lstset{
  language=Dart,
  basicstyle=\ttfamily\footnotesize,
  keywordstyle=\color{kw}\bfseries,
  keywordstyle=[2]\color{type},
  commentstyle=\color{cmt}\itshape,
  stringstyle=\color{str},
  numbers=left,
  numberstyle=\tiny\color{numgray},
  numbersep=10pt,
  showstringspaces=false,
  breaklines=true,
  breakatwhitespace=false,
  tabsize=2,
  frame=single,
  rulecolor=\color{codeframe},
  backgroundcolor=\color{codebg},
  framexleftmargin=6pt,
  xleftmargin=16pt,
  aboveskip=10pt,
  belowskip=6pt,
  literate=
    {á}{{\'a}}1 {é}{{\'e}}1 {í}{{\'i}}1 {ó}{{\'o}}1 {ú}{{\'u}}1
    {ñ}{{\~n}}1 {Á}{{\'A}}1 {É}{{\'E}}1 {Í}{{\'I}}1 {Ó}{{\'O}}1
    {Ú}{{\'U}}1 {Ñ}{{\~N}}1 {¡}{{\textexclamdown}}1 {¿}{{\textquestiondown}}1
}

% ---- Cajas de aviso ---------------------------------------------
\newtcolorbox{wherebox}{
  breakable, enhanced, colback=accent!6, colframe=accent,
  boxrule=0.4pt, left=8pt, right=8pt, top=6pt, bottom=6pt,
  fonttitle=\bfseries\color{white}, coltitle=white,
  colbacktitle=accent, title={Donde insertarlo}
}
\newtcolorbox{whatbox}{
  breakable, enhanced, colback=brand!6, colframe=brand,
  boxrule=0.4pt, left=8pt, right=8pt, top=6pt, bottom=6pt,
  fonttitle=\bfseries\color{white}, coltitle=white,
  colbacktitle=brand, title={Que hace este bloque}
}
\newtcolorbox{tipbox}{
  breakable, enhanced, colback=gold!12, colframe=gold!80!black,
  boxrule=0.4pt, left=8pt, right=8pt, top=6pt, bottom=6pt,
  fonttitle=\bfseries\color{black}, coltitle=black,
  colbacktitle=gold!55, title={Pruebalo}
}

\setlist{leftmargin=1.3em, topsep=3pt, itemsep=2pt}
\newcommand{\file}[1]{\texttt{\textcolor{brandDark}{#1}}}

% ==================================================================
\begin{document}

% ---- PORTADA ----------------------------------------------------
\begin{titlepage}
\thispagestyle{empty}
\vspace*{2.5cm}
\begin{center}
{\color{brand}\rule{\textwidth}{2pt}}\\[10pt]
{\Huge\bfseries\color{brandDark} Ultime Team Manager}\\[10pt]
{\LARGE\color{accent} Login paso a paso}\\[6pt]
{\large Formulario, validacion, estado y navegacion protegida en Flutter}\\[10pt]
{\color{brand}\rule{\textwidth}{2pt}}
\end{center}

\vspace{2cm}
\begin{center}
\begin{minipage}{0.82\textwidth}
\large
\textbf{Stack:} Flutter, Dart, Riverpod, go\_router, shared\_preferences.\\[6pt]
\textbf{Objetivo:} construir un login solido, offline y escalable a online, en
\textbf{6 fases} incrementales. En cada fase la app sigue compilando.\\[6pt]
\textbf{Integrantes:} Sebastian Ramirez \; \textbullet\; Joaquin Parraud.
\end{minipage}
\end{center}

\vfill
\begin{center}\small\textcolor{midgray}%
{Genera el PDF con \texttt{tectonic login-paso-a-paso.tex}}\end{center}
\end{titlepage}

\tableofcontents
\cleardoublepage

% ==================================================================
\chapter*{Introduccion}
\addcontentsline{toc}{chapter}{Introduccion}
\markboth{Introduccion}{}
% ==================================================================

Vamos a construir la pantalla de login de \textbf{Ultime Team Manager}. No es un
formulario suelto: es un login \emph{de verdad}, con arquitectura por capas, estado
gestionado con Riverpod, sesion persistente y navegacion protegida. Hoy funciona
\textbf{offline}; manana, cambiando una sola clase, funcionara contra Supabase.

\section*{El flujo que vamos a montar}
\begin{enumerate}
  \item La app arranca y \textbf{restaura} la sesion guardada (si la hay).
  \item Si no hay sesion, el \texttt{go\_router} nos lleva a \file{LoginScreen}.
  \item El usuario escribe correo y contrasena. El formulario \textbf{valida} los campos.
  \item Al pulsar \emph{Entrar}, el \texttt{AuthController} llama al repositorio,
        muestra un \textbf{spinner} y guarda la sesion si las credenciales son correctas.
  \item El router detecta el cambio de estado y \textbf{redirige} a \file{HomeScreen}.
  \item Desde Home, \emph{Cerrar sesion} borra los datos y vuelve al login.
\end{enumerate}

\section*{Las 6 fases}
\begin{enumerate}
  \item \textbf{Dependencias y arranque} --- anadimos paquetes y activamos Riverpod.
  \item \textbf{El dominio} --- la entidad \file{AppUser} y el contrato \file{AuthRepository}.
  \item \textbf{Repositorio local} --- login offline con sesion persistente.
  \item \textbf{Estado con Riverpod} --- el \file{AuthController} y sus providers.
  \item \textbf{La pantalla de login} --- formulario, validacion y estados de carga/error.
  \item \textbf{Navegacion protegida} --- \texttt{go\_router} con guardia de sesion.
\end{enumerate}

\begin{tipbox}
\textbf{Credenciales de demo} (modo offline): correo \texttt{demo@ultime.com},
contrasena \texttt{123456}. Las definiremos en la Fase 3 y podras cambiarlas.
\end{tipbox}

\begin{whatbox}
\textbf{Nota sobre el nombre del paquete.} El proyecto base se llama
\texttt{contador\_app}, por eso todos los \texttt{import} usan
\texttt{package:contador\_app/...}. Si mas adelante renombras el paquete, actualiza
esos imports de golpe.
\end{whatbox}

% ==================================================================
\chapter{Dependencias y arranque}
% ==================================================================
Antes de escribir logica, dejamos el proyecto listo: anadimos las dependencias y
envolvemos la app en un \texttt{ProviderScope} (el contenedor global de estado de
Riverpod).

\section{Anadir las dependencias}
\begin{wherebox}
Archivo \file{pubspec.yaml}. Dentro de la seccion \texttt{dependencies:}, justo
\textbf{debajo} de la linea \texttt{cupertino\_icons: \^{}1.0.8}.
\end{wherebox}

\begin{lstlisting}[language={}]
dependencies:
  flutter:
    sdk: flutter
  cupertino_icons: ^1.0.8
  flutter_riverpod: ^2.6.1      # gestion de estado
  go_router: ^14.6.2            # navegacion declarativa
  shared_preferences: ^2.3.3    # persistencia de la sesion
\end{lstlisting}

\begin{whatbox}
\begin{itemize}
  \item \texttt{flutter\_riverpod}: providers y \emph{notifiers} para el estado de autenticacion.
  \item \texttt{go\_router}: rutas por URL y \emph{redirect} (la guardia de sesion).
  \item \texttt{shared\_preferences}: guarda la sesion en el dispositivo (clave-valor).
\end{itemize}
\end{whatbox}

Ahora descarga los paquetes desde la terminal:
\begin{lstlisting}[language={},numbers=none]
flutter pub get
\end{lstlisting}

\section{Activar Riverpod en el arranque}
\begin{wherebox}
Archivo \file{lib/main.dart}. Sustituye el \texttt{import} y la funcion
\texttt{main()} actuales por lo siguiente (dejamos el resto del archivo igual por
ahora).
\end{wherebox}

\begin{lstlisting}
import 'package:flutter/material.dart';
import 'package:flutter_riverpod/flutter_riverpod.dart';
import 'package:contador_app/presentation/screens/counter/counter_screen.dart';

// Envolvemos la app en ProviderScope: sin el, ningun provider funciona.
void main() => runApp(const ProviderScope(child: MyApp()));
\end{lstlisting}

\begin{whatbox}
\texttt{ProviderScope} es la raiz donde Riverpod guarda el estado de todos los
providers. Debe envolver a \emph{toda} la app. De momento \file{MyApp} sigue
mostrando el contador; lo cambiaremos en la Fase 5.
\end{whatbox}

\begin{tipbox}
Ejecuta \texttt{flutter run}. La app debe seguir mostrando el contador sin errores.
Si compila, la base de Riverpod ya esta puesta.
\end{tipbox}

% ==================================================================
\chapter{El dominio: entidad y contrato}
% ==================================================================
La capa de dominio no sabe nada de Flutter ni de como se guardan los datos. Define
\textbf{que} es un usuario y \textbf{que operaciones} de autenticacion existen. Asi,
cambiar de ``offline'' a ``Supabase'' no toca ni la UI ni el dominio.

\section{La entidad AppUser}
\begin{wherebox}
Crea el archivo \file{lib/domain/entities/app\_user.dart} con este contenido.
\end{wherebox}

\begin{lstlisting}
// Representa al usuario autenticado dentro de la app.
class AppUser {
  final String id;
  final String email;

  const AppUser({required this.id, required this.email});
}
\end{lstlisting}

\begin{whatbox}
Un objeto inmutable (\texttt{final} + constructor \texttt{const}) con lo minimo que
necesitamos del usuario. Nada de contrasenas aqui: la entidad viaja por la UI.
\end{whatbox}

\section{El contrato AuthRepository}
\begin{wherebox}
Crea el archivo \file{lib/domain/repositories/auth\_repository.dart}.
\end{wherebox}

\begin{lstlisting}
import 'package:contador_app/domain/entities/app_user.dart';

// Contrato de autenticacion. La UI depende de esta abstraccion,
// no de una implementacion concreta (hoy local, manana Supabase).
abstract class AuthRepository {
  // Devuelve el usuario si hay sesion guardada, o null.
  Future<AppUser?> currentUser();

  // Inicia sesion. Lanza AuthException si las credenciales fallan.
  Future<AppUser> signIn({required String email, required String password});

  // Cierra la sesion y borra los datos guardados.
  Future<void> signOut();
}

// Error de dominio para credenciales incorrectas u otros fallos de auth.
class AuthException implements Exception {
  final String message;
  const AuthException(this.message);

  @override
  String toString() => message;
}
\end{lstlisting}

\begin{whatbox}
\begin{itemize}
  \item \texttt{AuthRepository} es \textbf{abstracto}: solo declara \emph{que} se puede
        hacer, no \emph{como}. En la Fase 3 escribimos la version offline.
  \item \texttt{currentUser()} permite restaurar la sesion al abrir la app.
  \item \texttt{AuthException} nos deja distinguir un error de credenciales de un
        \emph{bug}, y mostrar un mensaje claro al usuario.
\end{itemize}
\end{whatbox}

% ==================================================================
\chapter{Repositorio local con sesion persistente}
% ==================================================================
Aqui implementamos el contrato para el modo \textbf{offline}: validamos contra un
usuario de demo y guardamos la sesion con \texttt{shared\_preferences}.

\section{AuthRepositoryLocal}
\begin{wherebox}
Crea el archivo \file{lib/data/repositories/auth\_repository\_local.dart}.
\end{wherebox}

\begin{lstlisting}
import 'package:shared_preferences/shared_preferences.dart';
import 'package:contador_app/domain/entities/app_user.dart';
import 'package:contador_app/domain/repositories/auth_repository.dart';

// Implementacion OFFLINE del contrato de autenticacion.
class AuthRepositoryLocal implements AuthRepository {
  // Claves con las que guardamos la sesion en el dispositivo.
  static const _kUserId = 'auth_user_id';
  static const _kUserEmail = 'auth_user_email';

  // Usuario de demo para el modo offline. Cambialo si quieres.
  static const _demoEmail = 'demo@ultime.com';
  static const _demoPassword = '123456';

  @override
  Future<AppUser?> currentUser() async {
    final prefs = await SharedPreferences.getInstance();
    final id = prefs.getString(_kUserId);
    final email = prefs.getString(_kUserEmail);
    if (id == null || email == null) return null;
    return AppUser(id: id, email: email);
  }

  @override
  Future<AppUser> signIn({
    required String email,
    required String password,
  }) async {
    // Simula la latencia de una peticion de red real.
    await Future.delayed(const Duration(milliseconds: 600));

    final normalized = email.trim().toLowerCase();
    if (normalized != _demoEmail || password != _demoPassword) {
      throw const AuthException('Correo o contrasena incorrectos');
    }

    final user = AppUser(id: 'u1', email: normalized);
    final prefs = await SharedPreferences.getInstance();
    await prefs.setString(_kUserId, user.id);
    await prefs.setString(_kUserEmail, user.email);
    return user;
  }

  @override
  Future<void> signOut() async {
    final prefs = await SharedPreferences.getInstance();
    await prefs.remove(_kUserId);
    await prefs.remove(_kUserEmail);
  }
}
\end{lstlisting}

\begin{whatbox}
\begin{itemize}
  \item \texttt{signIn()} normaliza el correo, simula un retardo (para que el spinner
        se vea), y compara con las credenciales de demo. Si fallan, lanza
        \texttt{AuthException}; si aciertan, \textbf{guarda} la sesion.
  \item \texttt{currentUser()} reconstruye el \texttt{AppUser} desde lo guardado.
  \item \texttt{signOut()} borra las claves: sesion cerrada.
\end{itemize}
\end{whatbox}

\begin{tipbox}
\textbf{Escalable a online:} cuando quieras Supabase, creas
\texttt{AuthRepositorySupabase implements AuthRepository} y cambias \emph{una linea}
en el provider de la Fase 4. Ni la UI ni el dominio se enteran.
\end{tipbox}

% ==================================================================
\chapter{Estado con Riverpod}
% ==================================================================
El \file{AuthController} es el cerebro del login: guarda el estado (cargando, error,
usuario) y expone las acciones \texttt{signIn} y \texttt{signOut}. La UI solo
\emph{observa} este estado.

\section{Providers y estado de autenticacion}
\begin{wherebox}
Crea el archivo \file{lib/presentation/providers/auth\_provider.dart}.
\end{wherebox}

\begin{lstlisting}
import 'package:flutter_riverpod/flutter_riverpod.dart';
import 'package:contador_app/domain/entities/app_user.dart';
import 'package:contador_app/domain/repositories/auth_repository.dart';
import 'package:contador_app/data/repositories/auth_repository_local.dart';

// 1) Provider del repositorio. Para pasar a online, cambia SOLO esta linea
//    por AuthRepositorySupabase().
final authRepositoryProvider = Provider<AuthRepository>((ref) {
  return AuthRepositoryLocal();
});

// 2) Estados posibles de la sesion.
enum AuthStatus { unknown, authenticated, unauthenticated }

// 3) El estado inmutable que observa la UI.
class AuthState {
  final AuthStatus status;
  final AppUser? user;
  final bool isSubmitting;   // true mientras se procesa el login
  final String? errorMessage;

  const AuthState({
    this.status = AuthStatus.unknown,
    this.user,
    this.isSubmitting = false,
    this.errorMessage,
  });

  AuthState copyWith({
    AuthStatus? status,
    AppUser? user,
    bool? isSubmitting,
    String? errorMessage,
  }) {
    return AuthState(
      status: status ?? this.status,
      user: user ?? this.user,
      isSubmitting: isSubmitting ?? this.isSubmitting,
      errorMessage: errorMessage, // se limpia si no se pasa
    );
  }
}

// 4) El controlador: mantiene el AuthState y expone las acciones.
final authControllerProvider =
    NotifierProvider<AuthController, AuthState>(AuthController.new);

class AuthController extends Notifier<AuthState> {
  @override
  AuthState build() {
    _restoreSession();          // restaura la sesion al arrancar
    return const AuthState();   // estado inicial: unknown
  }

  AuthRepository get _repo => ref.read(authRepositoryProvider);

  Future<void> _restoreSession() async {
    final user = await _repo.currentUser();
    state = state.copyWith(
      status: user != null
          ? AuthStatus.authenticated
          : AuthStatus.unauthenticated,
      user: user,
    );
  }

  // Devuelve true si el login fue correcto.
  Future<bool> signIn(String email, String password) async {
    state = state.copyWith(isSubmitting: true); // enciende el spinner
    try {
      final user = await _repo.signIn(email: email, password: password);
      state = state.copyWith(
        status: AuthStatus.authenticated,
        user: user,
        isSubmitting: false,
      );
      return true;
    } on AuthException catch (e) {
      state = state.copyWith(isSubmitting: false, errorMessage: e.message);
      return false;
    }
  }

  Future<void> signOut() async {
    await _repo.signOut();
    state = const AuthState(status: AuthStatus.unauthenticated);
  }
}
\end{lstlisting}

\begin{whatbox}
\begin{itemize}
  \item \textbf{(1)} \texttt{authRepositoryProvider} es el \emph{unico} punto que
        conoce la implementacion concreta. Cambiarlo = cambiar de backend.
  \item \textbf{(3)} \texttt{AuthState} agrupa todo lo que la UI necesita pintar:
        estado, usuario, si esta enviando y el mensaje de error. Es inmutable;
        \texttt{copyWith} crea copias con cambios puntuales.
  \item \textbf{(4)} \texttt{AuthController} enciende \texttt{isSubmitting} antes de
        llamar al repositorio, y actualiza \texttt{state} segun el resultado.
        \texttt{build()} lanza la restauracion de sesion al crearse.
  \item \texttt{errorMessage} se limpia en cada \texttt{copyWith} que no lo pase: es
        un error \emph{transitorio}, no queremos que se quede pegado.
\end{itemize}
\end{whatbox}

% ==================================================================
\chapter{La pantalla de login}
% ==================================================================
Toca la UI: un formulario con validacion, un boton que muestra spinner mientras
carga y un mensaje de error si las credenciales fallan. Tambien creamos una Home
sencilla como destino tras entrar.

\section{LoginScreen}
\begin{wherebox}
Crea el archivo \file{lib/presentation/screens/auth/login\_screen.dart}.
\end{wherebox}

\begin{lstlisting}
import 'package:flutter/material.dart';
import 'package:flutter_riverpod/flutter_riverpod.dart';
import 'package:contador_app/presentation/providers/auth_provider.dart';

// Usamos ConsumerStatefulWidget: necesitamos estado local (los controllers
// de texto) y ademas leer/observar providers de Riverpod.
class LoginScreen extends ConsumerStatefulWidget {
  const LoginScreen({super.key});

  @override
  ConsumerState<LoginScreen> createState() => _LoginScreenState();
}

class _LoginScreenState extends ConsumerState<LoginScreen> {
  final _formKey = GlobalKey<FormState>();
  final _emailCtrl = TextEditingController();
  final _passCtrl = TextEditingController();
  bool _obscure = true;

  @override
  void dispose() {
    _emailCtrl.dispose();
    _passCtrl.dispose();
    super.dispose();
  }

  String? _validateEmail(String? value) {
    final text = value?.trim() ?? '';
    if (text.isEmpty) return 'Introduce tu correo';
    final regex = RegExp(r'^[\w.\-]+@[\w\-]+\.[\w.\-]+$');
    if (!regex.hasMatch(text)) return 'Correo no valido';
    return null;
  }

  String? _validatePassword(String? value) {
    final text = value ?? '';
    if (text.isEmpty) return 'Introduce tu contrasena';
    if (text.length < 6) return 'Minimo 6 caracteres';
    return null;
  }

  Future<void> _submit() async {
    FocusScope.of(context).unfocus();        // oculta el teclado
    if (!_formKey.currentState!.validate()) return;

    final ok = await ref
        .read(authControllerProvider.notifier)
        .signIn(_emailCtrl.text, _passCtrl.text);

    if (!ok && mounted) {
      final msg = ref.read(authControllerProvider).errorMessage ?? 'Error';
      ScaffoldMessenger.of(context)
          .showSnackBar(SnackBar(content: Text(msg)));
    }
    // Si ok == true, el router (Fase 6) redirige solo a Home.
  }

  @override
  Widget build(BuildContext context) {
    // Observamos solo isSubmitting para redibujar el boton/campos.
    final isSubmitting =
        ref.watch(authControllerProvider.select((s) => s.isSubmitting));

    return Scaffold(
      body: SafeArea(
        child: Center(
          child: SingleChildScrollView(
            padding: const EdgeInsets.all(24),
            child: Form(
              key: _formKey,
              child: Column(
                mainAxisSize: MainAxisSize.min,
                crossAxisAlignment: CrossAxisAlignment.stretch,
                children: [
                  const Icon(Icons.sports_soccer, size: 72, color: Colors.green),
                  const SizedBox(height: 16),
                  Text(
                    'Ultime Team Manager',
                    textAlign: TextAlign.center,
                    style: Theme.of(context).textTheme.headlineSmall,
                  ),
                  const SizedBox(height: 32),
                  TextFormField(
                    controller: _emailCtrl,
                    keyboardType: TextInputType.emailAddress,
                    textInputAction: TextInputAction.next,
                    enabled: !isSubmitting,
                    decoration: const InputDecoration(
                      labelText: 'Correo',
                      prefixIcon: Icon(Icons.email_outlined),
                      border: OutlineInputBorder(),
                    ),
                    validator: _validateEmail,
                  ),
                  const SizedBox(height: 16),
                  TextFormField(
                    controller: _passCtrl,
                    obscureText: _obscure,
                    enabled: !isSubmitting,
                    textInputAction: TextInputAction.done,
                    onFieldSubmitted: (_) => _submit(),
                    decoration: InputDecoration(
                      labelText: 'Contrasena',
                      prefixIcon: const Icon(Icons.lock_outline),
                      border: const OutlineInputBorder(),
                      suffixIcon: IconButton(
                        icon: Icon(_obscure
                            ? Icons.visibility_off
                            : Icons.visibility),
                        onPressed: () => setState(() => _obscure = !_obscure),
                      ),
                    ),
                    validator: _validatePassword,
                  ),
                  const SizedBox(height: 24),
                  FilledButton(
                    onPressed: isSubmitting ? null : _submit,
                    child: isSubmitting
                        ? const SizedBox(
                            height: 20,
                            width: 20,
                            child: CircularProgressIndicator(strokeWidth: 2),
                          )
                        : const Text('Entrar'),
                  ),
                  const SizedBox(height: 12),
                  Text(
                    'Demo: demo@ultime.com / 123456',
                    textAlign: TextAlign.center,
                    style: Theme.of(context).textTheme.bodySmall,
                  ),
                ],
              ),
            ),
          ),
        ),
      ),
    );
  }
}
\end{lstlisting}

\begin{whatbox}
\begin{itemize}
  \item \texttt{Form} + \texttt{\_formKey}: \texttt{validate()} ejecuta todos los
        \texttt{validator} y pinta los errores bajo cada campo.
  \item \texttt{\_validateEmail/\_validatePassword}: reglas de validacion (formato de
        correo y minimo 6 caracteres).
  \item \texttt{isSubmitting} deshabilita los campos y convierte el boton en un
        spinner: \textbf{estado de carga}.
  \item En error mostramos un \texttt{SnackBar}; \texttt{dispose()} libera los
        controllers para no tener fugas de memoria.
\end{itemize}
\end{whatbox}

\section{HomeScreen (destino tras el login)}
\begin{wherebox}
Crea el archivo \file{lib/presentation/screens/home/home\_screen.dart}.
\end{wherebox}

\begin{lstlisting}
import 'package:flutter/material.dart';
import 'package:flutter_riverpod/flutter_riverpod.dart';
import 'package:contador_app/presentation/providers/auth_provider.dart';

class HomeScreen extends ConsumerWidget {
  const HomeScreen({super.key});

  @override
  Widget build(BuildContext context, WidgetRef ref) {
    final user = ref.watch(authControllerProvider).user;

    return Scaffold(
      appBar: AppBar(
        title: const Text('Inicio'),
        actions: [
          IconButton(
            icon: const Icon(Icons.logout),
            onPressed: () =>
                ref.read(authControllerProvider.notifier).signOut(),
          ),
        ],
      ),
      body: Center(
        child: Text('Bienvenido, ${user?.email ?? 'jugador'}!'),
      ),
    );
  }
}
\end{lstlisting}

\begin{whatbox}
\texttt{ConsumerWidget} porque solo necesita \emph{leer} providers. El boton de la
\texttt{AppBar} llama a \texttt{signOut()}: al cambiar el estado, el router nos
devolvera al login.
\end{whatbox}

\section{Ver el login en marcha (temporal)}
\begin{wherebox}
Archivo \file{lib/main.dart}. Cambia el \texttt{import} del contador por el del
login, y pon \file{LoginScreen} como \texttt{home}. Es temporal: en la Fase 6 lo
sustituye el router.
\end{wherebox}

\begin{lstlisting}
// Sustituye esta linea:
import 'package:contador_app/presentation/screens/counter/counter_screen.dart';
// por esta:
import 'package:contador_app/presentation/screens/auth/login_screen.dart';

// ...y dentro de MaterialApp cambia:
//   home: const CounterScreen(),
// por:
      home: const LoginScreen(),
\end{lstlisting}

\begin{tipbox}
Ejecuta \texttt{flutter run}. Prueba a dejar campos vacios o un correo mal escrito:
deben salir los errores de validacion. Con credenciales incorrectas, veras el
\texttt{SnackBar} rojo. Con las de demo el login \emph{funciona}, pero aun no navega
a Home: eso lo arregla la Fase 6.
\end{tipbox}

% ==================================================================
\chapter{Navegacion protegida con go\_router}
% ==================================================================
Ultima pieza: el router decide a que pantalla ir segun el estado de la sesion. Si no
hay sesion, al login; si la hay, a Home. Y reacciona \textbf{automaticamente} cuando
el estado cambia.

\section{El router con guardia de sesion}
\begin{wherebox}
Crea el archivo \file{lib/config/router/app\_router.dart}.
\end{wherebox}

\begin{lstlisting}
import 'package:flutter/foundation.dart';
import 'package:flutter_riverpod/flutter_riverpod.dart';
import 'package:go_router/go_router.dart';
import 'package:contador_app/presentation/providers/auth_provider.dart';
import 'package:contador_app/presentation/screens/auth/login_screen.dart';
import 'package:contador_app/presentation/screens/home/home_screen.dart';

final routerProvider = Provider<GoRouter>((ref) {
  // Un notifier que "despierta" al router cuando cambia el estado de sesion.
  final refresh = ValueNotifier<AuthStatus>(AuthStatus.unknown);
  ref.listen(
    authControllerProvider.select((s) => s.status),
    (previous, next) => refresh.value = next,
    fireImmediately: true,
  );
  ref.onDispose(refresh.dispose);

  return GoRouter(
    initialLocation: '/login',
    refreshListenable: refresh,
    redirect: (context, state) {
      final status = ref.read(authControllerProvider).status;

      // Mientras restauramos la sesion, no decidimos nada todavia.
      if (status == AuthStatus.unknown) return null;

      final loggingIn = state.matchedLocation == '/login';
      final loggedIn = status == AuthStatus.authenticated;

      if (!loggedIn) return loggingIn ? null : '/login';
      if (loggingIn) return '/home';
      return null;
    },
    routes: [
      GoRoute(path: '/login', builder: (_, __) => const LoginScreen()),
      GoRoute(path: '/home', builder: (_, __) => const HomeScreen()),
    ],
  );
});
\end{lstlisting}

\begin{whatbox}
\begin{itemize}
  \item \texttt{redirect} es la \textbf{guardia}: si no hay sesion te manda a
        \texttt{/login}; si ya entraste, te saca del login a \texttt{/home}.
  \item Con \texttt{status == unknown} devolvemos \texttt{null}: esperamos a que
        termine la restauracion antes de redirigir (evita parpadeos).
  \item \texttt{refreshListenable} + \texttt{ref.listen}: cuando el
        \texttt{AuthController} cambia el estado, el \texttt{ValueNotifier} avisa al
        router y este vuelve a evaluar el \texttt{redirect}. Asi el login y el
        logout navegan solos.
\end{itemize}
\end{whatbox}

\section{Conectar el router en main.dart}
\begin{wherebox}
Archivo \file{lib/main.dart}. Reemplaza \textbf{todo} el contenido por esta version
final.
\end{wherebox}

\begin{lstlisting}
import 'package:flutter/material.dart';
import 'package:flutter_riverpod/flutter_riverpod.dart';
import 'package:contador_app/config/router/app_router.dart';

void main() => runApp(const ProviderScope(child: MyApp()));

// Ahora es ConsumerWidget para poder leer el routerProvider.
class MyApp extends ConsumerWidget {
  const MyApp({super.key});

  @override
  Widget build(BuildContext context, WidgetRef ref) {
    final router = ref.watch(routerProvider);
    return MaterialApp.router(
      debugShowCheckedModeBanner: false,
      theme: ThemeData(useMaterial3: true, colorSchemeSeed: Colors.green),
      routerConfig: router,
    );
  }
}
\end{lstlisting}

\begin{whatbox}
Pasamos de \texttt{MaterialApp} con \texttt{home:} a \texttt{MaterialApp.router} con
\texttt{routerConfig:}. \file{MyApp} ahora es \texttt{ConsumerWidget} para poder
\texttt{watch} del \texttt{routerProvider}. Ya no importamos ninguna pantalla
directamente: manda el router.
\end{whatbox}

\begin{tipbox}
\textbf{Prueba completa:}
\begin{enumerate}
  \item Arranca: te lleva al login (no hay sesion).
  \item Entra con \texttt{demo@ultime.com} / \texttt{123456}: navega a Home solo.
  \item \textbf{Cierra la app y vuelve a abrirla}: entra directo a Home (sesion persistida).
  \item Pulsa \emph{Cerrar sesion}: vuelves al login.
\end{enumerate}
Si todo esto ocurre, tienes un login completo, persistente y escalable a online.
\end{tipbox}

% ==================================================================
\chapter*{Resumen de archivos}
\addcontentsline{toc}{chapter}{Resumen de archivos}
\markboth{Resumen de archivos}{}
% ==================================================================
\begin{itemize}
  \item \file{pubspec.yaml} --- 3 dependencias nuevas \emph{(modificado)}.
  \item \file{lib/main.dart} --- ProviderScope y \texttt{MaterialApp.router} \emph{(modificado)}.
  \item \file{lib/domain/entities/app\_user.dart} --- entidad \file{AppUser}.
  \item \file{lib/domain/repositories/auth\_repository.dart} --- contrato + \file{AuthException}.
  \item \file{lib/data/repositories/auth\_repository\_local.dart} --- login offline.
  \item \file{lib/presentation/providers/auth\_provider.dart} --- providers y \file{AuthController}.
  \item \file{lib/presentation/screens/auth/login\_screen.dart} --- formulario y validacion.
  \item \file{lib/presentation/screens/home/home\_screen.dart} --- destino tras el login.
  \item \file{lib/config/router/app\_router.dart} --- navegacion protegida.
\end{itemize}

\vspace{6pt}
\noindent\textcolor{midgray}{\rule{\textwidth}{0.4pt}}\\[2pt]
\noindent{\small\textcolor{midgray}{Ultime Team Manager --- Login paso a paso ---
Sebastian Ramirez y Joaquin Parraud.}}

\end{document}
